\chapter{Introduction}

\label{introduction}

Gender bias in Machine Translation (MT) is a field of research that has gained popularity over the years, with a growing body of literature on issues of gender in MT \citep{savoldi2021gender}. For instance, in 2018, examples of gender bias in Google Translate started to appear and be shared on different social media platforms. These images showed how Google's MT engine systematically reproduced gender stereotypes when translating from Turkish, a genderless language, into English, a notional gender language, e.g. ``o bir mühendis'' became ``he is an engineer'', while ``o bir hemşire'' became ``she is a nurse'' (\cite{olson2018algorithm}). Studies on the topic can be mainly divided into two categories: The assessment of bias in commercial and academic systems, and debiasing methods that try to remove bias from MT outputs \citep{savoldi2021gender}. 

A common problem in this research field is that the term \textit{bias} is rarely discussed and/or defined in publications. The term itself is polysemous \citep{crawford2017keynote} and is used differently across disciplines. Defining bias is not trivial because it implies a normative process for which something, in this case MT output, is deemed as harmful to certain groups of people \citep{blodgett-etal-2020-language}. A literature review reveals how gender bias is often described as %the reason for 
the reproduction of gender stereotypes, whilst masculine forms disproportionally occur in MT outputs \citep{savoldi2021gender}. This application of the concept is influenced by assumptions on gender. For example, a survey of gender bias in the field of Natural Language Processing (NLP) demonstrates that gender is usually defined, explicitly or implicitly, as binary \citep{stanczak2021survey}. This is similar to MT research, where very few publications address non-binary genders (e.g. \cite{dev-etal-2021-harms,burtscher-2022-genderfairmt}).

The visibility of non-binary people has been increasing in the last few years. Numerous TV series, such as \textit{One Day at a Time} and \textit{Sex Education}, promote gender diversity. Celebrities like Sam Smith and Demi Lovato have come out as non-binary. Accordingly, the debate on gender-fair language (GFL) has attracted more and more attention, and  strategies that were previously considered inclusive enough are now outdated. GFL usually subsumes two main approaches, i.e. gender-neutral and gender-inclusive \citep{sczesny2016can,hornscheidt2021wie}. Gender-neutral means that sentences are usually structured to avoid gendered elements, for example by using gender-neutral words, passive constructions and indefinite pronouns. Gender-inclusive describes the use of linguistic means, e.g. neopronouns (\textit{hen} in Swedish), neomorphemes (\textit{e} in Spanish), typographical characters (gender star (*) in German) and/or symbols (\textit{ə} in Italian), to make all genders visible. 

GFL poses great challenges for humans and MT alike. Strategies vary greatly among languages and cultures due to their grammatical structures, i.e. how they express and encode gender, and the dominant ideologies and beliefs in a society (see e.g. \cite{lardellitranslating}). In countries like Germany and Austria, where feminist discourse has a strong tradition,  strategies have been, at least in part, proposed and used since the 1960s \citep{sczesny2016can}. Nowadays, the discussion around linguistic inclusivity also centres on genders beyond the binary. As a consequence, an increasing number of approaches and strategies have been proposed (see e.g. \cite{hornscheidt2021wie}). In other countries, for example Italy, where sexist attitudes are still dominant to a certain extent \citep{murgia2021stai}, GFL is less actively used and often focuses on the representation of women. Nevertheless, non-binary language is gaining visibility, even though it is often harshly criticised \citep{cosinonschwa}. 

With conceptualisations of translation shifting from being considered a simple process of cross-linguistic transfer to a complex process also involving culture and ideology, feminist work started to permeate Translation Studies (TS) from the 1970s onwards \citep{von2016translation}. Recently, paradigms such as Queer Translation \citep{baer2017introduction} have started to emerge. This is an interdisciplinary field that examines the intersections between translation and queer theory, focusing on how translation practices can both reveal and challenge normative understandings of gender and sexuality. It explores how translations can either obscure or highlight queer identities and experiences, and how the act of translation itself can be viewed through a queer lens. Such a queer lens questions fixed identities and categories, problematising the conventional notion of a clear-cut, hierarchical relationship between an original text and its translation, which influenced translation theory. Instead, it views translation as a dynamic and reciprocal process, where meanings are fluid and continuously negotiated \citep{spurlin2014queering,baer2017introduction}. However, non-binary genders are still neglected in both TS and MT research. This is a highly problematic issue for various reasons.

First, misgendering, i.e. the act of referring to someone with gendered forms that do not correspond to their own gender identity, is a form of identity invalidation that can cause psychological distress and thus negatively impact the mental health of trans and non-binary (TNB) people \citep{mclemore2015experiences,mclemore2018minority}. Second, due to their recently gained visibility, TNB people are increasingly appearing in translations. Even with the best intentions, subtitled or dubbed TV series as well as media reports on non-binary individuals are brimming with incorrect and misgendering references \citep{lopez2022trans,attig2023call,lardellitranslating}. Language has an impact on the visibility of genders and contributes to the (re)production of social norms and expectations. Misgendering or mistranslating leads to identity omission, invalidation and/or even ridicule because people might not understand what non-binary means. Given that MT is used on a large scale with millions of people accessing it every day for work, leisure, travel or understanding social media content, biased outputs can exacerbate the erasure and discrimination of non-binary people \citep{dev-etal-2021-harms}.

Most research on gender bias in MT tackles it from a purely technical perspective \citep{savoldi2021gender}. This, however, ignores the social and (cross-)linguistic issues for the emergence of gender bias as well as their interrelation with technical aspects. A merely technical approach fails to account for the social impact of technology use, which exacerbates the oppression of marginalised groups \citep{bender2021dangers,o2017weapons}.\footnote{The term \textit{marginalised} is used to highlight that discrimination and exclusion occur because of unequal power distributions in society.} A more fruitful exchange between TS and the MT community would be beneficial because joint research endeavours would reveal a much more complex picture of gender bias in both human translation and MT. 

The present work investigates the translation and post-editing (PE) of GFL beyond the binary. Twelve professional translators were recruited and asked either to translate or to post-edit machine translations of three brief English-language texts into German. The study was conducted online to aim at experimental realism, i.e. mirroring a real-life scenario in which participants could work in their usual environment. Participants completed a pre-study questionnaire with questions concerning socio-demographic information and their use of and previous experience with GFL. They were instructed to translate or post-edit each text using different  strategies, namely (i) gender-neutral rewording, (ii) gender-inclusive characters, and (iii) neosystems.\footnote{Noesystems are new gender systems based on neomorphemes} Process data were collected by combining non-participant observation, screen recordings, and retrospective interviews. The translations and post-edited texts were also analysed. This multi-method approach provides insights into the cognitive processes involved in gender-fair translation and PE. 

Similar studies in Translation Process Research (TPR) are usually conducted in a laboratory setting with strictly controlled variables \citep{carl2015post,jia2019does,nitzke2019problem}. This was not the case with the present study. The main objective was to gain first insights into the translation and PE process but, more importantly, to put participants in a position to experiment with  and gain insights into their personal experiences and points of view. The cross-linguistic expertise of translators is often neglected when tackling gender issues in MT. However, such expertise is fundamental to understanding the (cross-)linguistic implications of GFL use and investigating the feasibility and acceptability of different strategies. For this reason, while translation and PE times, as well as screen activity tracking data, provide interesting results, the core of the analysis is represented by the interviews during which participants expressed their opinions on the tasks and commented on the use of technology. 

\textcite{attig2023call} calls for a community-informed translation practice where non-binary people are involved in the decision-making process concerning the selection of GFL strategies. This participatory model emphasises the importance of engaging with queer communities to ensure translations reflect authentic language practices, as demonstrated in the analysis of Spanish and French translations of non-binary pronouns in the sitcom \textit{One Day at a Time}.\footnote{\url{https://en.wikipedia.org/wiki/One_Day_at_a_Time_(2017_TV_series)}}

Building upon this framework, the concept of \textit{co-creation} offers a complementary perspective that deepens the collaborative nature of translation practices. For instance, \textcite{perez2013co} explores how digital media and technological advancements have shifted the nature of subtitling, creating new collaborative and transformative practices. The article argues that fansubbing and other participatory subtitling activities demonstrate a move away from traditional, industry-controlled subtitling and towards a co-creative model where users actively participate in the creation and adaptation of media content. 

Co-creation therefore involves stakeholders as equal partners throughout the design and decision-making processes, fostering a shared vision and iterative collaboration. In the context of TS, this approach aligns with participatory action research methodologies that include experiential experts, such as queer and non-binary individuals, in the translation process. \textcite{gromann2023participatory} illustrate this through a case study focusing on German MT, highlighting the importance of context dependency and the desire for customizable solutions to avoid identity invalidation.

Integrating co-creation into community-informed translation practices ensures that non-binary individuals are not merely consulted but are active collaborators in the translation process. This approach acknowledges their expertise and lived experiences, leading to translations that are more inclusive and representative. By embracing co-creation, TS can move toward more equitable practices that honour the diversity of gender identities.

Although the involvement of non-binary people was marginal -- as a matter of fact, only one non-binary translator volunteered to take part in the study -- this work attempts to promote a change in research practices when tackling issues of gender bias and non-binary representation by involving a group whose expertise and role in circulating ideas are fundamental at a societal level, i.e. professional translators. 

%This is a question of power imbalances and people in privileged positions should not further silence marginalised communities but rather support their right to linguistic self-determination. For this reason, I advocate for community-informed and intersectional translation practices that take into account the different needs and sensitivities of the queer community and linguists, also considering the needs of other marginalised groups such as people with disabilities.

The main findings from this study show that neosystems were the most difficult strategy to use, with a considerable increase in screen activity tracking, perceived cognitive effort, and the number of mistakes made when they were applied. In terms of time, PE was found to be faster than translation, but no statistically significant difference emerged among assignments. Finally, participants clearly indicated a preference for a combination of gender-neutral rewording and gender-inclusive characters. The present work provides insights into the strengths and limitations of different strategies as well as into particularly challenging gender-specific terms, which can be of inspiration for MT debiasing and further research on gender-fair MT.

The present work is structured as follows: 
\begin{itemize}
\setlength\itemsep{-0.5em}
    \item Chapter \ref{introduction} presents the research questions, assumptions, motivation, and objectives underlying this research.
    \item Chapter \ref{chapter2-1} delves into the theoretical framework of the present work, which is highly multidisciplinary. First, this monograph is set in the context of feminist and queer TS, which are the result of shifts in the conception of translation from a mere linguistic transfer to a complex process of cross-linguistic, cross-cultural, ideological and political negotiation. Second, relevant concepts from gender studies, such as sex, gender, and gender identity, are introduced. The objective of this chapter is to create awareness that gender is a biopsychosocial construct and that there are numerous gender identities beyond man and woman. Finally, gender as a linguistic feature is introduced and the interplay between language, culture, and society is described. 
    \item Chapter \ref{Chapter2-2} is dedicated to gender issues in MT: Gender bias and its sociotechnical implications are discussed. Subsequently, the second section of the chapter focuses on PE as an increasingly common practice in the translation industry.
    \item Chapter \ref{Chapter2-3} is dedicated to an overview of linguistic strategies in English, German, and Italian. Language-independent good practices for gender fairness are also briefly mentioned. Finally, related work from TS and MT research, as well as work that combines both disciplines is summarised.
    \item Chapter \ref{methodology} details the research methodology. It introduces the multi-method approach used for the investigation of the research object and focuses on the experiment preparation, data collection, and analysis.
    \item Chapter \ref{results} provides an overview of the main findings. It briefly describes the participants' profiles, then concentrates on the analysis of translation and PE times, screen activity tracking data, participants' experience, the use of strategies during the experiment, and the bias in the machine translations post-edited in the study.
    \item Chapter \ref{Discussion} discusses the results. First, the different  approaches, i.e. gender-neutral rewording, gender-inclusive characters and neosystems, and the translation and PE process are compared. Second, the implications for language professionals and MT research are addressed. Finally, the author's positionality and the limitations of this study are described.
    \item Chapter \ref{Conclusion} presents a concise overview and a concluding analysis of the present work.  
\end{itemize}
Following the introduction of the book topic, the research problem, and the structural outline of the work, the following section presents the research questions and the primary assumptions underpinning this study.


\section{Research questions and main assumptions}
This study represents the first endeavour to investigate GFL beyond the binary in human translation and MT. It is both a process and product-oriented study partially inspired by work in TPR. The following research questions were devised:
\begin{itemize}
\setlength\itemsep{-0.5em}
    \item \textbf{Q1}: How do translators apply different GFL strategies in translation and PE and what are the cognitive processes involved?
    \item \textbf{Q2}: To what extent is the integration of GFL in translation and PE successful?
    \item \textbf{Q3}: What are the differences in effort between gender-fair translation and PE?
\end{itemize}
The first question can also be divided into three sub-questions, namely:
\begin{itemize}
\setlength\itemsep{-0.5em}
    \item How do professional translators perceive the task of using GFL? What are the differences between strategies and what are the advantages and disadvantages of each?
    \item To what extent do translators’ subjective impressions of and experience with GFL correlate with other indicators of cognitive processes, i.e. translation times, and screen activity tracking data?
    \item Which gender-specific terms are particularly challenging in the translation and PE process?
\end{itemize}

The main assumptions concerning the findings of this study are that language professionals use GFL differently, e.g. they select different neutralisation strategies or gender-inclusive characters. In terms of cognitive processes, neosystems are expected to be more challenging than the other strategies because they are largely unknown and actively change and disrupt traditional grammar rules. It is also expected that translators make more mistakes with them than when using gender-neutral rewording and gender-inclusive characters. Longer translation times should correspond to a higher perceived cognitive effort, and PE should be faster than translation. Different gender-specific terms are expected to be challenging based on the strategy applied, e.g. personal pronouns and professions when using gender-neutral rewording, and pronouns and articles when using gender-inclusive characters. Finally, terms specific to a gender, e.g. \textit{Nanny} in German, are expected to cause greater difficulties in the translation process.

The present work is also influenced by more general assumptions that affected its realisation. Gender identity terms used in this book reflect Western ideologies, current views of gender, and an Anglocentric bias in research and academia. Identity terms and views on gender are different in non-Western languages and cultures \citep{lester2018trans}; these are not addressed in the present work. I also acknowledge that the views reflected in this work are not only the result of a social construction but they also change over time. Labels and language used here may become outdated in the future.

%I briefly discuss my privilege because I conduct research that involves and perhaps affects a marginalised group\footnote{The term marginalised is used to highlight that discrimination and exclusion are perpetrated because of unequal power distribution among people in society. It is preferred over minorities since discriminated groups, for instance, women are not a minority.} to which I do not belong. For this purpose, the concept of ``intersectionality'' needs to be introduced. The term was coined by \textcite{crenshaw1989demarginalizing} and describes how race, class, gender, and other individual characteristics overlap with one another. An example to understand how the concept works is the lawsuit by Emma DeGraffenreid against General Motors (GM) in 1976. GM was accused of discriminating against black people but won the lawsuit. When the judges verified that the company hired black people and white women, they believed there was no discrimination against black women, ignoring that black women face discrimination precisely because they are both black and women (\cite[142f]{crenshaw1989demarginalizing}). One analogy that can be used to understand intersectionality is that of crossroads. Every road represents a line of identity, for example, gender, sexuality, race, and class. The unique combination (or intersection) of these roads describes the multiple oppression faced by a single person. Fig. \ref{fig:intersectionality} offers an intuitive instrument to describe the different axes of one's identity. I am an Italian, white, able-bodied, cisgender, middle-class, researcher at a public institution in Austria. I never experienced forms of oppression and discrimination like people from marginalised groups. This also means that I have blind spots of which I am not always aware and which influence my beliefs and my research as well as the interpretation of my research results.

%\begin{figure}[ht]    
%\centering
%\includegraphics[width=0.8\textwidth]{figures/Intersectionality_wheel_of_power_and_privilege_from_University_of_British_Columbia.png}
  %  \caption{The wheel of power and privilege}
   % \label{fig:intersectionality}
%\end{figure}

Lastly, I also recognise that GFL is an ideologically/politically charged and often controversial topic. Since one of its objectives is to undermine or overturn the status quo by making marginalised groups visible and prominent in language, GFL use tends to cause strong reactions of outcry and mockery \citep{vergoossen2020four,cosinonschwa}. While people who support GFL use tend to be accused of promoting a specific political and ideological agenda, its opposers are not. Therefore, I refuse to discuss non-binary representation in language from an ideological point of view, concentrating instead on challenges of a (cross-)linguistic nature. I maintain that using gendered forms which correspond to a person's gender identity is a matter of respect and not of ideology. Nevertheless, I will briefly touch upon my positionality as a researcher in Chapter \ref{Discussion}.

\section{Research relevance and objectives}
The present work is highly relevant to both translation studies and practice as well as MT. Research on gender bias often neglects non-binary individuals. This represents a considerable research gap that this study tries to address. Queer TS predominantly focuses on sexuality, e.g. the translation of gayspeak, or on trans people who identify with either the male or the female gender \citep{baer2017introduction}. This is a significant limitation from a theoretical and empirical perspective. %explain why - discourse around visibility

From the point of view of the translation industry, audiovisual translation continues to grow as a field and non-binary characters are increasingly depicted in TV series that are dubbed and subtitled in many languages. In the next few years, it can be expected that translation professionals will be increasingly confronted with non-binary language and thus need to be able to reproduce the meaning of the source text without erasing gender identities beyond the binary. Unfortunately, translations are still brimming with misgendering \citep{lopez2022trans,attig2023call,lardellitranslating}. As language technologies are increasingly integrated into translation workflows, language professionals should be aware of the biased outputs produced by MT and be able to adjust them.

Research on gender bias in NLP and MT often adopts an implicit binary conception of gender. While some studies acknowledge that gender is not binary, gender issues are tackled from a binary perspective due to the lack of groundwork \citep{sun2019mitigating,stanczak2021survey}. TS and NLP should collaborate on gender issues that are common in human and MT. Integrating GFL in MT while also taking into account the multiplicity of the available strategies is not possible without the perspective of language professionals who have a greater awareness of (cross-)linguistic issues resulting from the adoption of specific  strategies.

Finally, the present work is also relevant on a broader societal level. Greater awareness of gender and linguistic issues as well as the use of GFL can further promote gender equality \citep{sczesny2016can}. Dismantling old-fashioned gender norms, stereotypes, and expectations may have a positive impact on the lives of many people who daily experience forms of discrimination and oppression because they do not strictly adhere to said norms, stereotypes, and expectations (\cite{barker2019life}). Academia and language professionals can become ambassadors for linguistic and societal change.

Having considered the research relevance, the present work pursues the following objectives: (i) gain insights into the translation and PE process concerning different GFL strategies, (ii) provide a methodological framework and materials to conduct further studies in other language pairs, (iii) create and share novel datasets with non-binary language that can be used for further research endeavours, (iv) contribute to the exchange on gender bias between TS and computational linguistics/MT, and (v) produce a tentative best practice model for language professionals, which can also provide inspiration for debiasing techniques in language technologies.

\section{Philosophical stance and research approach}
\textcite[10]{saldanha2014research} advocate for critically examining one’s worldview before embarking on a research project. This is fundamental in the selection of appropriate research methods, the validation of the findings, and the (self-)understanding of the researcher’s motivation. Considerations of what knowledge is and how it is produced are therefore essential and relate to two complementary positions, the ontological and the epistemological. 

Ontology is the branch of philosophy concerned with the nature and relations of being. In social research, ontology pertains to “the way the social world is seen to be and what can be assumed about the nature and reality of the social phenomena that make up the social world” \citep[23]{matthews2010research}. Epistemology, on the other hand, concerns the theory of knowledge acquisition and production \citep[23]{matthews2010research}. Three different ontological positions may be identified, i.e. objectivism, constructivism, and realism, which are broadly associated with three corresponding epistemological positions, i.e. positivism, interpretivism, and realism.

Objectivism posits that social phenomena exist independently of social actors \citep[33]{bryman2016social}. This position is usually linked to positivism, which conceives science as value-free, objective, based on quantitative data and rigorous hypothesis-testing \citep[28]{bryman2016social}. Constructivism ``implies that social phenomena and categories are not only produced through social interaction but that they are in a constant state of revision'' \citep[33]{bryman2016social}. It is logically connected to interpretivism, which ``prioritises people’s subjective interpretations and understandings of social phenomena'' \citep[28]{matthews2010research}. The constructivist-interpretivist position acknowledges that research is a culturally and socially situated practice, privileging qualitative data and understanding (\cite[30]{bryman2016social}). Finally, realism, both as an ontological and epistemological position, recognises that social phenomena can have a reality independent of social actors, but the scientist’s conceptualisation is just one way of knowing them (\cite[11]{saldanha2014research}, \cite[29]{bryman2016social}).

Finally, two basic orientations to research can be adopted, namely inductive and deductive. The inductive approach, often associated with qualitative research, involves deriving general principles from specific instances or observations. This method allows for the exploration of patterns and the generation of theories based on empirical data, offering a bottom-up perspective \citep[14-15]{saldanha2014research}. Conversely, the deductive approach is characterised by a top-down methodology where hypotheses are formulated based on existing theories and subsequently tested through empirical observation. This approach is typically linked to quantitative research and aims to validate or refute theoretical constructs, providing a structured framework for testing specific translation hypotheses \citep[14--15]{saldanha2014research}.

The present work and the author’s personal beliefs are influenced by feminist theory, which asserts that knowledge and truth are situated in a specific societal, cultural, and historical context that inevitably shapes and influences them \citep{haraway1998science}. This is in line with the social-constructivist stance and the conceptualisation of translation as a socially and culturally embedded practice. However, the methodological approach employed in this study is mixed methods. Even though quantitative and qualitative methods are often seen as mutually exclusive, their distinction is not always clear-cut and their combination might be fruitful to understanding the research object (\cite[5]{luker2009salsa}, \cite[37]{bryman2016social}). Mixed methods make it possible to analyse the translation process, the translation product and the social context, which in turn allows for a triangulation of data from different perspectives. The first perspective is the objective one and relates to the time needed to complete translation and PE assignments as well as the activities carried out, such as pauses and searches. The second is more subjective and it relates to the personal experience of language professionals. The analysis of data, for instance the coding process, followed an inductive approach. The result is a robust work that provides numerous insights into GFL translation and PE in the language pair English-German.

This chapter outlined the topic of the present monograph, namely the translation and PE of GFL to promote respectful representation of non-binary individuals. It also emphasised the underresearched nature of the topic in both TS and computational linguistics, underscoring the relevance of this work. Following an overview of the structural outline, the research questions, objectives, and philosophical stance underpinning this study were discussed. The theoretical background is presented in Chapters \ref{chapter2-1}, \ref{Chapter2-2}, and \ref{Chapter2-3}.


